\section{Two nonlinear boundary layers}
\label{sec:v4-boundary}
Fix one oriented channel and suppress its index and the smooth geometric
parameter. All bounds below include each fixed number of parameter
derivatives on the compact families specified in
Section~\ref{sec:geometric-results}. Shrinking the common contact box is
allowed once, independently of the flight number. The functions
$E_j=W_j-jg$ and $b_j=-W_{j,uv}/d_j^0$ refer to the actual nonlinear
stationary action and its normalized twist, not their quadratic models.

\subsection{The half-line construction}
For a starting type $b\in\{0,1\}$ put
$\sigma_i^{(b)}=\sqrt{c_{1-((b+i)\bmod2)}}$. The quadratic half-line
interior Hessian acts on sequences indexed by $i\ge1$, with fixed zero
boundary at $i=0$. Its diagonal entries are $2c_{(b+i)\bmod2}/g$
and its adjacent entries are $-1/g$. Its inverse is
\begin{equation}\label{eq:v4-half-green}
 G^{(b)}_{ik}=\frac{g\sigma_i^{(b)}\sigma_k^{(b)}}{2c\sinh\gamma}
 \left(e^{-\gamma|i-k|}-e^{-\gamma(i+k)}\right),\qquad i,k\ge1.
\end{equation}
The formula follows either from the recurrence and its unit jump or
from \eqref{eq:v3-green} by sending the right endpoint to infinity.
It is a bounded positive operator on $\ell^2(\mathbb N)$, and is
bounded on $\sup_i\rho^{-i}|x_i|$ whenever
$e^{-\gamma_-}<\rho<1$. Every fixed parameter derivative obeys the
same assertion after leaving a strict exponential margin.

\begin{lemma}[Semi-infinite stationary segments]\label{lem:v4-halfline}
For $|u|<u_*$ there is a unique small stationary half-line segment
$x^{(b)}(u)=(x_i^{(b)}(u))_{i\ge0}$ with $x_0^{(b)}=u$ and
$x_i^{(b)}\to0$. It is smooth, and on a common neighborhood
\begin{equation}\label{eq:v4-half-estimate}
 \begin{split}
 x_i^{(b)}(u)&=\frac{\sigma_i^{(b)}}{\sigma_0^{(b)}}
                      e^{-i\gamma}u+O(u^2\rho^i),\\
 |\partial_\xi^\alpha\partial_u^r x_i^{(b)}(u)|
                  &\le C_{\alpha,r}\rho^i
                     |u|^{\max(1-r,0)}.
 \end{split}
\end{equation}
The boundary action
\begin{equation}\label{eq:v4-boundary-action}
 S_b(u)=\sum_{i\ge0}
  \left[\ell_{(b+i)\bmod2}(x_i^{(b)}(u),x_{i+1}^{(b)}(u))-g\right]
\end{equation}
is smooth and strictly convex near zero. It satisfies
\begin{equation}\label{eq:v4-boundary-hessian}
 S_b(0)=S_b'(0)=0,\qquad
 a_b:=S_b''(0)=\frac{c\sinh\gamma}{g c_{1-b}}>0.
\end{equation}
These assertions are analytic on a common smaller neighborhood for
jointly analytic graph families.
\end{lemma}
\begin{proof}
Use the half-line Green operator in the weighted contraction from
Lemma~\ref{lem:g-relative}. The linear solution is the first term of
\eqref{eq:v4-half-estimate}. The nonlinear gradient is local, vanishes
quadratically, and its derivative has norm $O(|u|)$ on the weighted
ball of radius $C|u|$. Thus the map has contraction constant at most
one half on a common ball. Differentiating its fixed-point equation
leaves the same invertible operator on the highest derivative; every
other term is a product of a local derivative and weighted lower
ones. The convolution estimate used in \eqref{eq:v3-weighted-green}
proves \eqref{eq:v4-half-estimate} by induction.

Uniqueness also holds among all sufficiently small stationary sequences
tending to zero, not just among elements of that weighted ball. Their
difference solves a homogeneous Jacobi equation whose averaged Hessian
is strictly diagonally dominant. Its inverse on bounded sequences is
bounded by the reciprocal dominance margin: write the Hessian as its
positive diagonal times $I-T$, with $\|T\|_\infty<1$.
With zero left boundary the difference is zero. This argument applies
to every fixed small contact box.

Every summand of \eqref{eq:v4-boundary-action} is $O(u^2\rho^{2i})$.
Its differentiated sums converge by the same weighted bounds. Interior
stationarity cancels the first-variation terms, leaving only the left
endpoint derivative; the boundary term at a truncation tends to zero.
At quadratic order this derivative is
$(c_bu-x_1^{(b),\mathrm{lin}})/g$. Since
$c_b\sigma_0^{(b)}/\sigma_1^{(b)}=c$, it equals $a_bu$.
This proves \eqref{eq:v4-boundary-hessian}, hence strict convexity.
The uniform complex contraction and Cauchy estimates, as in
Lemma~\ref{lem:g-relative}, prove the analytic assertion.
\end{proof}

Let $H_b(u)$ be the half-line interior Hessian evaluated on this
segment, and write $\Delta H_b(u)=H_b(u)-H_b(0)$. It is tridiagonal.
Locality and \eqref{eq:v4-half-estimate} imply
\begin{equation}\label{eq:v4-half-trace}
 \|\Delta H_b(u)\|_1\le C|u|,
 \qquad
 \|\partial_\xi^\alpha\partial_u^r\Delta H_b(u)\|_1
                      \le C_{\alpha,r}.
\end{equation}
For both finite and half-line operators, the accompanying bounds are
$\|\partial_\xi^\alpha G\|_{\mathrm{op}}\le C_\alpha$ and
$\|\partial_\xi^\alpha\partial_u^r(G\Delta H)\|_1\le C_{\alpha,r}$.
The first follows by differentiating $HG=I$; the second uses the
trace-class factor $\Delta H$ and its derivatives. These are precisely
\eqref{eq:v5-operator-trace-bounds}, and apply also to compressed
operators in the comparisons below.
Here and below $\|\cdot\|_1$ on operators is the trace norm. Indeed,
the sum of absolute entries bounds that norm, and the entries are
bounded by sums of neighboring endpoint weights. Define
\begin{align}
 U_b(u)&=\sum_{i\ge0}\log\left(
 g[-\ell_{(b+i)\bmod2,uv}(x_i^{(b)},x_{i+1}^{(b)})]\right),
 \label{eq:v4-half-product}\\
 B_b(u)&=\exp\left\{U_b(u)
 -\sum_{m\ge1}\frac{(-1)^{m+1}}m
       \tr[(G^{(b)}\Delta H_b(u))^m]\right\}.
 \label{eq:v4-half-amplitude}
\end{align}
This also defines the Fredholm determinant in the denominator of
$B_b$: no additional determinant normalization is implicit. The series
converges absolutely after requiring
$\|G^{(b)}\Delta H_b\|<1/2$. Each trace is bounded by one trace
norm times $m-1$ operator norms. The edge sum converges by endpoint
localization. Consequently $B_b>0$, $B_b(0)=1$, and every fixed
mixed derivative is uniformly bounded. Differentiation of the trace
series retains a trace-class factor, so its bound does not involve the
number of sites of a truncation.

\subsection{Relative factorization}
\begin{theorem}[Nonlinear two-boundary factorization]
\label{thm:v4-factorization}
There are a common endpoint box and $\tau\in(0,1)$ such that for every
fixed $k$,
\begin{equation}\label{eq:v4-factorization}
 \|E_j-S_0\oplus S_p\|_{C^k}
 +\|b_j-B_0\otimes B_p\|_{C^k}\le C_k\tau^j,
 \qquad p=j\bmod2.
\end{equation}
The norms include mixed geometric-parameter and endpoint derivatives;
$S_0\oplus S_p$ means $S_0(u)+S_p(v)$ and
$B_0\otimes B_p$ means $B_0(u)B_p(v)$. The action difference has
zero value and first endpoint differential at zero, and therefore is
$O(\tau^j(|u|+|v|)^2)$. The amplitude difference vanishes at zero.
The fixed-order constants, but not $j$, may depend on higher boundary
norms. Both limits are analytic for analytic families.
\end{theorem}
\begin{proof}
We give the gluing and determinant estimates separately. Reverse a
half-line segment starting at type $p$ at the right endpoint. Its
$i$th coordinate in a length-$j$ indexing is $x_{j-i}^{(p)}(v)$;
the reversal is consistent because the flight action is symmetric
under reversal together with interchange of endpoint types. For
$0<i<j$ put
\[
 z_i=x_i^{(0)}(u)+x_{j-i}^{(p)}(v),\qquad z_0=u,\quad z_j=v.
\]
The endpoint overwrites cost $O(\rho^j)$ in every fixed differentiated
norm. In each interior stationarity equation the linear terms cancel.
The remaining interaction is a sum of products of left and right
coordinates or their neighboring values. Its size is bounded by
$C\rho^i\rho^{j-i}=C\rho^j$. Thus, for the interior residual $r(z)$,
\begin{equation}\label{eq:v4-gluing-residual}
 \|\partial^\alpha r(z)\|_{\ell^1}
                          \le C_\alpha(1+j)\rho^j.
\end{equation}
The assertion for derivatives follows from the local chain rule;
parameter derivatives of the separated linear exponentials introduce
only fixed powers of indices, absorbed by a strict exponential margin.

Let $y$ be the finite stationary segment already constructed in
Lemma~\ref{lem:g-relative}. The identity
\[
 r(z)-r(y)=\left[\int_0^1 H_{\mathrm{int}}(y+t(z-y))\,\dd t\right](z-y)
\]
has a symmetric strictly diagonally dominant matrix with a uniform
positive margin. Its inverse has bounded $\ell^\infty$ norm by diagonal
splitting and a geometric series; symmetry gives the same $\ell^1$
bound. It follows that $\|y-z\|_{\ell^1}\le C(1+j)\rho^j$.
Differentiating this identity leaves the same uniformly invertible
operator on the highest difference derivative. All other terms contain
lower difference derivatives and uniformly bounded local coefficients.
Induction proves the corresponding $\ell^1$ estimate at every fixed
order, with a polynomial in $j$ multiplying an exponential. Fixing a
slower exponential absorbs every such fixed polynomial.

For the action, subtract the two half-line sums. The missing tails are
exponentially small. On overlapping edges the difference between the
value at the sum and the sum of the two values minus $g$ is a local
interaction bounded by $C\rho^j$. The replacement of $z$ by $y$ is
controlled by its $\ell^1$ error and bounded local derivatives.
The same argument after differentiation proves the first term of
\eqref{eq:v4-factorization}. Both finite and limiting actions have
zero value and gradient at zero; Taylor's integral remainder yields
the additional quadratic vanishing.

For the normalized twist, the finite cofactor formula gives
\begin{equation}\label{eq:v4-finite-log}
 \log b_j=
 \sum_{i=0}^{j-1}\log(g[-\ell_{i\bmod2,uv}(y_i,y_{i+1})])
 -\log\det(I+G_j\Delta H_j),
\end{equation}
where $G_j=(H^0_{\mathrm{int}})^{-1}$ and
$\Delta H_j=H_{\mathrm{int}}-H^0_{\mathrm{int}}$.
The edge sum converges exponentially to $U_0(u)+U_p(v)$ by exactly
the preceding local interaction and tail estimates. It remains to
control the determinant relatively.

Here are the needed finite-dimensional-to-half-line details. For
$j\ge6$ choose $r=\lfloor j/3\rfloor$ and retain the first and last
$r$ interior sites. Set to zero all perturbation entries outside their
two diagonal blocks, including crossing entries. The trace-norm error
is bounded by
\begin{equation}\label{eq:v4-det-tail}
 C_k(1+j)^{M_k}\rho^r
\end{equation}
after $k$ fixed derivatives. This uses tridiagonality and the decay of
the two endpoint layers, not a dimension times operator-norm estimate.
Within each retained block, replacing the finite segment by its own
half-line segment has an exponentially small entrywise sum, by
\eqref{eq:v4-gluing-residual} and the opposite tail.

For a perturbation supported in these blocks, the identity
$\det(I+AB)=\det(I+BA)$ reduces the determinant to that support.
The left compression of $G_j$ differs from the compression of
$G^{(0)}$ by an exponentially small operator norm. The reversed right
compression differs in the same way from $G^{(p)}$. To check the claim,
use \eqref{eq:v3-green}: after expanding hyperbolic sines into
exponentials, the reflection from the remote boundary has a factor
$e^{-\gamma(2j-i-k)}$ for $i,k\le r$. The two off-diagonal block
kernels have a factor $e^{-\gamma(j-2r)}$. Summing at most $r^2$
entries and differentiating any fixed number of times only adds a
polynomial factor. Consequently the compressed Green matrix differs
exponentially from the direct sum of the two compressed half-line
operators. The same estimate holds after multiplication by the bounded
trace-norm perturbation.

For completeness, if $T,T'$ are the resulting trace-class operators
with norms at most $q<1$, telescoping powers gives
\[
 |\tr(T^m-(T')^m)|\le m q^{m-1}\|T-T'\|_1.
\]
Summation with coefficient $1/m$ proves local Lipschitz continuity of
the logarithmic determinant in trace norm. Its differentiated series
has the same property with a fixed polynomial factor in $m$, still
summable. This justifies each truncation and comparison above, for all
fixed mixed derivatives. Sending the block size to infinity at either
end costs $O(\rho^r)$ again. The direct-sum logarithm is the sum of
the half-line logarithms in \eqref{eq:v4-half-amplitude}.

We have proved exponential convergence of \eqref{eq:v4-finite-log}
to $\log B_0(u)+\log B_p(v)$. Exponentiation preserves the bounds,
since all logarithms range over one bounded interval. Choose $\tau<1$
slower than the finitely many exponential rates in the gluing,
truncation, and off-diagonal estimates. Polynomial factors at each
fixed derivative order are absorbed in $C_k$. The finitely many
$j<6$ are absorbed by enlarging those constants. At zero the finite
and limiting normalized twists are one, which proves the last
vanishing assertion. Uniform complex versions of the same estimates
prove analyticity. In particular the argument never divides an
absolute action error by the exponentially small reference twist.
\end{proof}

\subsection{Physical limiting laws}
For $p\in\{0,1\}$ let $q_p=\sqrt{a_0a_p}$ and define, for $d>0$,
\begin{equation}\label{eq:v4-limit-integral}
 \mathcal F_p(d)=\frac{q_p}{\pi d^2}
 \int_{\mathbb R^2}(d-S_0(u)-S_p(v))_+B_0(u)B_p(v)\,\dd u\dd v.
\end{equation}
The integral is restricted to the common contact box; for the collars
below its support lies strictly inside that box, so the displayed
extension by zero is unambiguous.

\begin{theorem}[Fixed-offset long-bridge law]\label{thm:v4-law}
There is $d_*>0$ independent of $j$ such that $\mathcal F_p$ has a
smooth right extension to $[0,d_*]$, with $\mathcal F_p(0)=1$ and
$\mathcal F_p>0$. On this interval,
\begin{equation}\label{eq:v4-F-convergence}
 \left\|\frac{2A\sinh(j\gamma)}{d^2}\Pr(E_{j,e}(d))
                   -\mathcal F_p(d)\right\|_{C^k}
 \le C_k\tau^j,\qquad p=j\bmod2,
\end{equation}
for every fixed mixed parameter and offset derivative order.
The conditional endpoint and first-residual-time law converges to
\begin{equation}\label{eq:v4-limit-law}
 \frac{B_0(u)B_p(v)}{I_p(d)}
 \one_{\{r>0,\ S_0(u)+S_p(v)+r<d\}}\,\dd u\dd v\dd r,
 \quad I_p(d)=\int(d-S_0-S_p)_+B_0B_p\,\dd u\dd v.
\end{equation}
For the scaled variables $(u/\sqrt d,v/\sqrt d,r/d)$ the convergence
has bounded-Lipschitz error at most $C\tau^j$, uniformly for
$0<d\le d_*$. Any fixed number of collisions at the two ends may be
included, with their limiting coordinates given by the corresponding
half-line segments. The two ends are coupled by the common residual
energy constraint; \eqref{eq:v4-limit-law} is not a product probability.
\end{theorem}
\begin{proof}
At parity $p$, \eqref{eq:g-ratio} and
\eqref{eq:v4-boundary-hessian} give $d_j^0=q_p/\sinh(j\gamma)$.
Full-phase integration therefore gives exactly
\[
 F_j(d)=\frac{q_p}{\pi d^2}\int(d-E_j)_+b_j\,\dd u\dd v.
\]
To compare this integral down to $d=0$, use the constructive Morse
maps of Lemma~\ref{lem:g-radial} for both $E_j$ and $S_0\oplus S_p$.
Their quadratic matrices are uniformly positive. The matrix
$2\int_0^1(1-s)D^2E(sy)\dd s$, its positive square root and the inverse
coordinate map depend locally Lipschitz-continuously in each fixed
$C^k$ norm on $E$, at the cost of finitely many additional derivatives.
This follows by differentiating the defining inverse equation and
using the uniform inverse differential; the positive-square-root
map has bounded derivatives on the common compact positive cone.
Thus their inverse maps and transformed amplitudes differ by
$C_k\tau^j$. The inverse maps both vanish at zero, so their difference
at $w$ is also $O(\tau^j|w|)$.

Put $w=\sqrt{2d}\,z$. The normalized integrals are over the same unit
disk with weight $1-|z|^2$. Their odd Taylor terms vanish. Taylor's
integral remainder, to two additional endpoint orders for each requested
offset derivative, bounds the difference of the smooth functions of
$d$ by $C_k\tau^j$. This proves \eqref{eq:v4-F-convergence} including
its derivatives at zero; it is not differentiation of a moving sharp
boundary without justification. The quadratic integral at zero is
$\pi d^2/q_p$, proving $\mathcal F_p(0)=1$. Positivity and a smaller
common collar give a uniform positive lower bound.

For conditional laws retain $r=d\eta$ and the fixed domain
$\mathcal U$ in \eqref{eq:v3-latent-domain}. The normalized densities
there differ in $L^1$ by $C\tau^j$, since the denominators are uniformly
positive. The two endpoint maps, after division by $\sqrt d$, differ
uniformly by $C\tau^j$. Couple the common part of the latent densities
and use these maps; the remaining mass is $O(\tau^j)$ on a bounded
scaled domain. This proves the bounded-Lipschitz estimate. A fixed
number of end coordinates follows from the differentiated gluing
estimate and its vanishing at zero. First variations give directions
from neighboring coordinates when desired. The explicit density
\eqref{eq:v4-limit-law} follows by returning to $(u,v,r)$ coordinates.
\end{proof}

\begin{corollary}[Nonlinear weights at a tied onset]\label{cor:v4-weights}
Suppose a fixed finite set of selected ground channels has the same
gap. Let $\gamma_{\min}$ be the least of their exponents. At a fixed
$0<d\le d_*$, along a parity subsequence their conditional weights in
the complete maximal-count event converge to
\[
 \frac{\one_{\{\gamma_e=\gamma_{\min}\}}\mathcal F_{e,p}(d)}
 {\sum_{e':\gamma_{e'}=\gamma_{\min}}\mathcal F_{e',p}(d)}.
\]
The convergence is exponential. At $d=0$ the minimizing leading
weights are equal, whereas at positive offset the boundary laws
supply their corrections.
\end{corollary}
\begin{proof}
Divide the finite sum in \eqref{eq:g-competition} by
$d^2e^{-j\gamma_{\min}}/A$. Theorem~\ref{thm:v4-law} and
$\csch(j\gamma)=2e^{-j\gamma}/(1-e^{-2j\gamma})$ give the limit.
All limiting denominators are positive. A finite positive gap between
distinct exponents and the factorization rate give an exponential
error; if all exponents tie only the latter rate is needed.
\end{proof}

\begin{corollary}[The first poles of the onset series]\label{cor:v4-poles}
For a fixed geometry put
$\mathcal G(z,d)=\sum_{j\ge1}z^j\Pr(E_{j,e}(d))/d^2$.
There is $\nu>0$ such that, with $x=ze^{-\gamma}$,
\begin{equation}\label{eq:v4-poles}
 \mathcal G(z,d)=
 \frac{\mathcal F_1(d)x+\mathcal F_0(d)x^2}{A(1-x^2)}
                      +\mathcal H(z,d),
\end{equation}
where $\mathcal H$ is holomorphic for $|z|<e^{\gamma+\nu}$ and
smooth in $d\in[0,d_*]$. Every fixed geometric and offset derivative of the remainder is
holomorphic on strictly smaller common domains. Derivatives of the
principal term are taken explicitly in \eqref{eq:v4-poles}; a parameter
derivative moving $\gamma$ can raise its pole order. The simple-pole
statement below concerns the undifferentiated series.
The positive singularity is a simple pole at $z=e^\gamma$, with
residue $-e^\gamma[\mathcal F_0(d)+\mathcal F_1(d)]/(2A)$.
The possible pole at $-e^\gamma$ is removable exactly when
$\mathcal F_0(d)=\mathcal F_1(d)$.
\end{corollary}
\begin{proof}
The coefficient in question is
\[
 \frac{e^{-j\gamma}}{A(1-e^{-2j\gamma})}F_j(d)
 =\frac{e^{-j\gamma}}A\mathcal F_{j\bmod2}(d)
                       +O(e^{-(\gamma+\nu)j}),
\]
for any positive $\nu$ strictly smaller than $-\log\tau$ and
$2\gamma$. Fixed parameter derivatives add only fixed powers of $j$,
absorbed by a smaller strict margin. The error series is normally
convergent, as are its stated derivatives. Summing its even and odd
leading terms gives \eqref{eq:v4-poles}. Positivity gives the nonzero
positive residue. At the negative point the numerator is
$\mathcal F_0-\mathcal F_1$, proving the removability criterion.
\end{proof}

The series in \eqref{eq:v4-poles} uses a different physical window
$jg+d$ at each coefficient. It is not the generating function of all
itineraries in one observation window. Its continuation is a consequence
of a nonlinear relative-flux limit for these specified events.
